\begin{summarybox}
	\textbf{\textcolor{CSummary}{一句话核心}}
	复合函数奇偶性由内函数主导：内偶则偶；内奇同外（奇偶性与外函数相同）。
	\medskip
	\textbf{\textcolor{CSummary}{使用条件}}
	\begin{itemize}[leftmargin=*, itemsep=0.5em, topsep=0.4em]
		\item \textbf{条件 1（定义域对称）：} 内函数$g(x)$的定义域必须关于原点对称。
		\item \textbf{条件 2（值域包含）：} 外函数$f(u)$的定义域必须包含内函数$g(x)$的值域。
	\end{itemize}
	\medskip
	\textbf{\textcolor{CSummary}{关键提醒}}
	\begin{itemize}[leftmargin=*, itemsep=0.5em, topsep=0.4em]
		\item \textbf{易错点（忽视前提）：} 应用口诀前先检查定义域对称性。
		\item \textbf{检查项（外层参与）：} 内奇时，务必考虑外层的奇偶性，不可直接断定整体为奇。
	\end{itemize}
\end{summarybox}
